\section*{Chapitre \ref{ch:m3:defi-credit-and-leverage} --- Crédit et levier dans la finance décentralisée}

\begin{solution}{exo:m3:defi-credit-and-leverage:1}
À 50~\%~: emprunter $4\% \times 50/90 = 2.22\%$, déposer $2.22\% \times 0.5 \times 0.9 = 1.00\%$. À 95~\%~: emprunter $4\% + 60\% \times 5/10 = 34\%$,
déposer $34\% \times 0.95 \times 0.9 = 29.07\%$.
\end{solution}

\begin{solution}{exo:m3:defi-credit-and-leverage:2}
$100\,000/(50 \times 0.825) = 2\,424.24$ dollars.
\end{solution}

\begin{solution}{exo:m3:defi-credit-and-leverage:3}
La transaction se termine soit avec le prêt et les frais remboursés, soit elle est annulée dans son intégralité, auquel cas les \omterm{def:m3:blockchains-for-traders:key}{jetons} n'ont jamais quitté le \omterm{def:m3:defi-credit-and-leverage:pool}{pool de prêt}. Il n'existe aucun état dans
lequel le prêt reste en cours.
\end{solution}

\begin{solution}{exo:m3:defi-credit-and-leverage:4}
$50\,000/2\,500 \times 1.05 = 21$ ether.
\end{solution}

\begin{solution}{exo:m3:defi-credit-and-leverage:5}
Non~: $(1 + \beta)\theta = 1.1 \times 0.825 = 0.9075$, au-dessus du \omterm{def:m3:defi-credit-and-leverage:ltv}{facteur de santé} de 0.80, si bien que chaque liquidation le fait baisser davantage
(\cref{prop:m3:defi-credit-and-leverage:liq})~; la position se dirige vers une dette irrécouvrable.
\end{solution}

\begin{solution}{exo:m3:defi-credit-and-leverage:6}
Déposer le \omterm{def:m3:defi-credit-and-leverage:staking}{jeton de staking liquide}, emprunter de l'ether contre lui, miser cet ether pour obtenir davantage de ce \omterm{def:m3:blockchains-for-traders:key}{jeton}, puis recommencer, et encaisser le \omterm{def:m3:defi-credit-and-leverage:staking}{rendement de staking} sur plusieurs
fois le capital, sous déduction du taux d'emprunt sur la dette. La stratégie perd si le taux d'emprunt dépasse le \omterm{def:m3:defi-credit-and-leverage:staking}{rendement de staking} (pics du taux d'utilisation), ou si le \omterm{def:m3:blockchains-for-traders:key}{jeton}
se négocie en dessous de l'ether qui le sous-tend et que les boucles sont liquidées.
\end{solution}

\begin{solution}{exo:m3:defi-credit-and-leverage:7}
Avec un gonflement de 1.93 si les \omterm{def:m3:blockchains-for-traders:key}{jetons} achetés deviennent sans valeur, de 1.72 s'ils conservent une valeur de USD~1~: emprunter 60~\% d'une valeur gonflée et abandonner un collatéral
valant la valeur réelle rapporte dès que $0.6k$ dépasse un augmenté du coût du gonflement par dollar de collatéral.
\end{solution}

\begin{solution}{exo:m3:defi-credit-and-leverage:8}
Surcollatéralisé au cours de l'\omterm{def:m3:blockchains-for-traders:contract}{oracle}. Si ce cours peut être déplacé, ou si le collatéral ne peut pas être vendu près de ce cours lors d'une chute, ou si les liquidateurs n'interviennent pas à temps,
le collatéral peut valoir moins que la dette~: la protection n'est aussi solide que l'\omterm{def:m3:blockchains-for-traders:contract}{oracle} et la liquidité du collatéral.
\end{solution}

\begin{solution}{pb:m3:defi-credit-and-leverage:1}
\textbf{1.} $0.6 \times 5 = $ USD~3 millions. \textbf{2.} $k^* = 50/(0.6 \times 5) = 16.7$. \textbf{3.} $2 \times (\sqrt{16.7} - 1) = $ USD~6.16
millions, en achetant $1 - 1/\sqrt{16.7} = 75.5\%$ des \omterm{def:m3:blockchains-for-traders:key}{jetons} du marché, soit 1.51 million. \textbf{4.} $50 - 5 - 6.16 = $ USD~38.84 millions~; USD~40.35
millions si les 1.51 million de \omterm{def:m3:blockchains-for-traders:key}{jetons} conservent une valeur de USD~1. \textbf{5.} Le pool est vide~: un gonflement supplémentaire ne fait qu'accroître les coûts. \textbf{6.} 1.93. \textbf{7.} Cette
attaque ne rapporte jamais~: faire monter un marché de USD~20 millions coûte plus que tout emprunt qu'elle permet. \textbf{8.} $k^*$ double pour atteindre 33.3, mais l'attaque reste
rentable à partir d'un gonflement de 4.98 et peut rapporter USD~35.4 millions~: diviser par deux le \omterm{def:m3:defi-credit-and-leverage:ltv}{ratio prêt/valeur} ne suffit pas. \textbf{9.} Un plafond d'environ USD~5.8 millions sur les emprunts
adossés à ce collatéral~: le gain $\min(B, 3k) - 5 - 2(\sqrt{k} - 1)$ (en millions de USD) n'est alors jamais positif. \textbf{10.} L'attaquant devrait
maintenir le prix gonflé pendant une demi-heure, ce qui multiplierait le coût et les pertes d'arbitrage tant que le prix reste hors d'alignement. \textbf{11.} Eisenberg a fait grimper le MNGO
sur les bourses alimentant l'\omterm{def:m3:blockchains-for-traders:contract}{oracle}, en faisant monter le prix publié de plus de treize fois en 30 minutes, et a utilisé ses positions gonflées comme collatéral pour
retirer plus de USD~110 millions. \textbf{12.} Sur une requête fondée sur la Rule 29, qui demande si les éléments de preuve pouvaient soutenir le verdict, le
juge a annulé les deux premiers chefs et prononcé un jugement d'acquittement sur le troisième (23 mai 2025). \textbf{13.} Les fournisseurs de liquidité du pool, au prorata, sauf si le
protocole dispose d'une réserve ou d'un backstop. \textbf{14.} Elle dure plus longtemps qu'une transaction et exige le capital propre de l'attaquant~; une attaque par \omterm{def:m3:defi-credit-and-leverage:flash}{flash loan} emprunte
le capital et doit s'achever en une seule transaction, si bien qu'elle ne peut exploiter que des prix qui bougent à l'intérieur de celle-ci. \textbf{15.} Geler le collatéral, abaisser son
\omterm{def:m3:defi-credit-and-leverage:ltv}{ratio prêt/valeur}, plafonner les emprunts adossés à ce collatéral, changer l'\omterm{def:m3:blockchains-for-traders:contract}{oracle} et voter sur les recouvrements. \textbf{16.} Les deux~: la conception a rendu l'attaque bon marché, et la question de savoir si la
conduite était illicite relève des tribunaux. \textbf{17.} Uniquement avec des plafonds d'emprunt calibrés sur le coût de manipulation des sources de son \omterm{def:m3:blockchains-for-traders:contract}{oracle}, ou pas du tout. \textbf{18.} Par le
collatéral qu'ils acceptent et les plafonds qui s'y appliquent, leurs \omterm{def:m3:blockchains-for-traders:contract}{oracles}, leur historique de dettes irrécouvrables et leurs réserves, ainsi que le taux d'utilisation (la capacité de retrait). \textbf{19.}
\emph{Résultat nommé~:} un gonflement de 16.7 fois épuise le pool~; il coûte USD~6.16 millions sur le marché source (plus les USD~5 millions de collatéral
abandonné) contre USD~50 millions empruntés, soit un gain d'environ USD~38.8 millions. \textbf{20.} Le seul élément qu'il connaît du monde, et donc ce
qu'achète un attaquant.
\end{solution}

\begin{solution}{iq:m3:defi-credit-and-leverage:1}
Collatéral pondéré par les seuils, rapporté à la dette. En dessous de un, n'importe qui peut rembourser une partie de la dette (dans la limite du facteur de clôture) et prendre du collatéral avec une décote, le
\omterm{def:m3:defi-credit-and-leverage:ltv}{bonus de liquidation}.
\iqlookfor{le ratio, le seuil et la liquidation ouverte à tous.}
\end{solution}

\begin{solution}{iq:m3:defi-credit-and-leverage:2}
Un prêt qui doit être remboursé dans la transaction, faute de quoi la transaction est annulée. Il exige une exécution atomique de code arbitraire entre contrats, ce
qu'une \omterm{def:m3:blockchains-for-traders:chain}{blockchain} assure et que le règlement traditionnel ne fournit pas.
\iqlookfor{l'atomicité tenant lieu de collatéral.}
\end{solution}

\begin{solution}{iq:m3:defi-credit-and-leverage:3}
Le pool n'a pas d'échéance et doit garder des \omterm{def:m3:blockchains-for-traders:key}{jetons} disponibles pour les retraits~; une pente forte au-delà du taux d'utilisation cible force le remboursement et attire des apports
avant que le pool ne soit entièrement prêté, tandis qu'une pente douce en dessous maintient des emprunts bon marché en temps normal.
\iqlookfor{la liquidité pour les retraits.}
\end{solution}

\begin{solution}{iq:m3:defi-credit-and-leverage:4}
À partir de la liquidité du collatéral (profondeur des marchés que lit l'\omterm{def:m3:blockchains-for-traders:contract}{oracle} et où vendent les liquidateurs), de sa volatilité sur la durée de la liquidation, du
coût de manipulation de son \omterm{def:m3:blockchains-for-traders:contract}{oracle}, et d'un plafond d'emprunt pour que la perte maximale soit bornée~; à réexaminer à mesure que les marchés évoluent.
\iqlookfor{horizon de liquidation, profondeur et coût de manipulation.}
\end{solution}

\begin{solution}{iq:m3:defi-credit-and-leverage:5}
Le rendement était versé et non gagné, et l'ancrage reposait sur la valeur du \omterm{def:m3:blockchains-for-traders:key}{jeton} flottant excédant les \omterm{def:m3:blockchains-for-traders:stable}{stablecoins} en circulation~; lorsque la demande s'est retournée, les rachats
ont créé des \omterm{def:m3:blockchains-for-traders:key}{jetons} flottants que personne ne voulait. Surveiller le ratio entre la valeur du \omterm{def:m3:blockchains-for-traders:key}{jeton} flottant et l'offre de \omterm{def:m3:blockchains-for-traders:stable}{stablecoins}, source du rendement, ainsi que les
flux de rachats.
\iqlookfor{réflexivité et subvention.}
\end{solution}

\begin{solution}{iq:m3:defi-credit-and-leverage:6}
Surveiller les \omterm{def:m3:defi-credit-and-leverage:ltv}{facteurs de santé} à partir de l'état de la chaîne et des mises à jour d'\omterm{def:m3:blockchains-for-traders:contract}{oracle} en attente~; simuler chaque liquidation avec le chemin de vente du collatéral et le \omterm{def:m3:blockchains-for-traders:gas}{gas}~; financer avec un
\omterm{def:m3:defi-credit-and-leverage:flash}{flash loan} et vendre dans la même transaction afin de ne conserver aucun inventaire~; n'enchérir sur une position via des bundles qu'à hauteur du profit attendu~; et plafonner l'exposition
au collatéral qui ne peut pas être vendu d'un seul coup.
\iqlookfor{financement et cession atomiques, et discipline d'enchère.}
\end{solution}
